Learning-based routers for wireless sensor networks are rarely compared with strong shortest-path baselines or with a lifetime bound, so their benefit is hard to judge. We benchmark multi-hop tree routing without aggregation against the Chang--Tassiulas fractional-flow bound, extended to an oracle for time-varying traffic, with every control bit charged: battery reports and route advertisements that encode every next hop. On 30 topologies under four traffic scenarios, battery-weighted Dijkstra reaches 85.2\% of the bound. Coordinating the routes chosen within a round adds about 4.5 points. Chang and Tassiulas's flow augmentation, run in within-round increments, reaches 89.9\%, and a sequential booked-load Dijkstra 89.7\%; at this operating point the two are indistinguishable (difference $-0.2$ points, 95\% CI $[-1.3, 0.8]$). An LP planner re-solved every round reaches 96.6\%, which shows that the bound is on average within 3.6\% of what a sequence of trees can achieve. On fresh topologies neither coordination mechanism dominates. Flow augmentation is ahead in small and clustered networks. Booked-load Dijkstra is ahead where traffic funnels hardest: with the sink at a corner it reaches 85.6\% against 79.3\%. Under lossy links with retransmissions and contention, coordination keeps its lifetime gain but lowers the delivery ratio. A reconstruction of a recent multi-agent reinforcement-learning router, with its reward and selection protocol corrected, collapses to direct transmission, although its own reward prefers Dijkstra routes eightfold; we report this as a reproducibility failure. When the sink can afford an LP per round, the planner is usually the better choice.
